\documentclass{article}
\usepackage{/globals/math}
\usepackage{/globals/de}
\usepackage{/globals/anki}
\ankiprefix{beweise}
\ankitopic{Beweise}
\ankishow
\begin{document}
\section{Beweise} 
Ein \emph{Beweis} ist eine Herleitung welche formal zeigt dass eine bestimmte Aussage korrekt sein muss, wenn davon ausgegangen wird dass andere Aussagen wahr sind.
\ankinote{def}{Definition}{Herleitung zeigt dass Aussage stimmt}

Die zu beweisene Aussage ist ein \emph{Theorem}, manchmal auch ein Satz \emph{genannt}. Es kann Hilfsätze (\emph{Lemma}s) und unmittelbar Folgende Aussagen (\emph{Korollar}s) geben.

\subsection{Beweismethoden}
Es gibt mehrere Methoden einen Beweis $A \implies B$ aufzustellen.

\begin{description}
 \item[Direkter Beweis] Mithilfe von ein bisschen Logik und Regeln wird direkt $A \implies B$ bewiesen. Dabei kann es auch zwischenschritte geben, \zb
\[
 A \implies X
 \dtext{und}
 A \implies Y
 \dtext{und}
 (X \land Y) \implies B
\]
\ankinote{direkter-beweis}{Direkter Beweis}{A zu B}
 \item[Beweis durch Kontraposition] Es wird die Kontraposition $\lnot B \implies \lnot A$ bewiesen. 
\ankinote{kontraposition}{Kontraposition}{\(\lnot B \implies \lnot A\)}
 \item[Beweis durch Widerspruch] Es wird davon ausgegangen dass $A \land \lnot B$ wahr ist und es wird gezeigt dass dies zu einem Widerspruch führt. Deshalb, solange bereits sicher ist dass $A$ wahr ist, kann $\lnot B$ nicht wahr sein, heißt $B$ muss wahr sein.
\ankinote{widerspruch}{Widerspruch}{Angenommen A und B führt zu Widerspruch}
 \item[Vollständige Fallunterscheidung] Es werden alle Möglichkeiten herausgearbeitet inwiefern $A$ wahr sein kann. Es wird gezeigt dass in all diesen Fällen $B$ wahr sein muss.
\ankinote{fallunterscheidung}{Fallunterscheidung}{Alle Fälle wie A wahr sein kann führen zu B}
 \item[Vollstände Induktion] Soll ein Prädikat hinter einem Allquantor bewiesen werden, wie ein $\forall x \in \N A(x)$ bewiesen werden, kann bewiesen werden dass $A(0)$ wahr ist, und dass $\forall x \in \N A(x) \implies A(x+1)$
\ankinote{induktion}{Induktion}{\(\forall x \in \N A(x) \implies A(x+1)\)}
\end{description}
\ankinote{methoden}{Beweismethoden}{Direkter Beweis, Kontraposition, Widerspruch, Fallunterscheidung, Induktion}
\end{document}